\section{유일하게 남은 후보: $SL_2(p)$ Weil 격자의 텐서족}\label{sec:leech}

자유 작용의 한계(정리 \ref{thm:fpf})를 넘으려면 ``고정점을 가진 원소는 모두 정칙형(경계)이고, 고정점 없는 순환부분군이 아주 많은'' 군이 필요하다. 3회차에서 살아남은 유일한 후보를 5·10·11·12회차에서 조사했다. 결론을 먼저 적는다.

\begin{center}\fbox{\begin{minipage}{0.93\textwidth}
\textbf{현황.} $p=5$ ($E_8$): 안정. $p=13$ (Leech): 정확히 경계 \tagPC. $(p+1)/2$ 합성수($p=17,29,41,\dots$): 불안정 \tagP. \textbf{$p=37$: 불안정} --- 72차원 짝 unimodular Weil 격자에 $N(\nu)=63<81$ 인 노름 6 벡터가 존재(독립 구현 3개 일치, 부동소수점 계산) \tagN. \textbf{$p=61$ (120차원)}: \S\ref{sec:p61} 참조. 성립하더라도 고정된 $p$ 는 상수 이득일 뿐이고, 차수 개선에는 무한히 많은 $p$ 와 아직 증명되지 않은 해석적 정리가 필요하다.
\end{minipage}}\end{center}

\subsection{후보족과 그것이 특별한 이유}
$p\equiv1\pmod4$ 소수, $G=SL_2(p)$, $W$ = 차원 $2(p-1)$ 의 충실한 유리 표현($W\otimes\C=2(\psi+\psi')$, $\psi,\psi'$ 는 차수 $(p-1)/2$ 의 홀 Weil 지표, 지표체 $F=\Q(\sqrt p)$, Schur 지수 2). $\Delta=\mathrm{End}_G(W)$ 는 $F$ 위의 정부호 사원수대수(두 실자리에서만 분기; $p\equiv5\pmod 8$ 이면 $(-1,-1)_F$), $r=(p-1)/4$, 이득 $g=\tfrac{p(p^2-1)}{16r^2}\approx p$. 이득의 대부분은 \emph{비분리 토러스}의 $\sim p^2/2$ 개 순환부분군(위수 $\approx p/2$, 고정점 없음)에서 온다.

\begin{remark}[왜 하필 $SL_2(p)$ 인가]
보조정리 \ref{lem:unimod} 에 의해, 안정(또는 경계)인 unimodular 격자에서는 소수 위수 $q$ 의 모든 원소 $g$ 가 (i) 자유롭게 작용하거나 ($\chi_W(g)=-\chi_W(1)/(q-1)$) (ii) $\chi_W(g)=0$ 이어야 한다. (i)이 성립하는 소수에서는 Sylow 부분군이 순환 또는 일반화 사원수군이고, (ii)는 전형적으로 $\chi$ 가 $q$-결함 0 일 때($\chi(1)_q=|G|_q$) 성립한다. 따라서 이득 $\lesssim|G|/\chi(1)^2$ 가 크려면 ``자유 소수'' 부분이 $|G|$ 의 대부분을 차지해야 한다. 모든 Sylow 부분군이 순환/사원수인 비가해군은 본질적으로 $SL_2(p)$ 뿐이고(Suzuki), $SL_2(p)$ 는 주기적 코호몰로지를 가지면서도 고정점 없는 선형 표현을 갖지 않는 고전적 예다. Weil 지표는 $p-1$ 의 홀수 소인수에서 결함 0, $p$ 와 2 에서 자유이며, 비분리 토러스의 소수 $q\mid\tfrac{p+1}2$ 에서 자유일 조건이 정확히 ``$(p+1)/2$ 가 소수'' 다. 같은 논법으로 $Sp_{2n}(p)$ ($n\ge2$)의 Weil 격자, $SL_2(p^f)$ ($f\ge2$), $p\equiv3\pmod4$ 의 경우, Suzuki 군 등은 배제된다(고정공간이 큰 원소의 지표값이 0 이 아님) \tagH\ (분류 정리가 아니라 조사 수준의 논증).
\end{remark}

\subsection{$p=13$: Leech 격자는 정확히 경계 안정}
\begin{theorem}\label{thm:leech}\tagPC\ (정수 연산; \texttt{rounds/R5\_leech13/}; 독립 심사자가 자체 코드로 재현)\
$g\in\mathrm{Aut}(\Lambda_{24})$ 를 위수 13의 자기동형, $\eta=\sum_{j=1}^{12}(\tfrac j{13})g^j$ 라 하자 ($\eta^2=13$).
\begin{enumerate}
\item 최소벡터 $w$ 에 대해 $s(w):=\tfrac12\,w\cdot\eta w\in\{0,\pm2,\pm4\}$ 이고 개수는 $78624,\ 45864\ (\times2),\ 13104\ (\times2)$; $s=\pm6$ 은 없다(순환 Gram 행렬의 양반정부호 조건을 유리수 연산으로 전수 확인 + 196560개 직접 계산).
\item $C_{Co_0}(g)=\langle g\rangle\times SL_2(3)$. $\Delta=\Q\langle SL_2(3)\rangle\otimes\Q(\eta)$ 의 order $\{x:x\Lambda_{24}\subset\Lambda_{24}\}$ 는 \emph{극대}(Hurwitz$\otimes O_F$ 위 지수 4)이고 class number 1. $\Lambda_{24}$ 는 그 위의 랭크 3 Hermite 격자이며 $\Delta$-선형 자기동형군은 위수 정확히 $2184$ 로 $SL_2(13)$ 과 동형.
\item 모든 $w\ne0$ 에 대해 $N_{F/\Q}(\nu(w))=\tfrac{13|w|^4}{16}-\tfrac{s(w)^2}4\ge9=r^2$, 모든 $\Delta$-직선 $U$ 에 대해 $\det(\Lambda_{24}\cap U)\ge3^8$ ($R_U\le1$), 모든 $\Delta$-평면에 대해 $R_U\le3^8/2^{16}=0.1001$.
\item 등호 $R_U=1$ 은 \textbf{정확히 $1638=3\times546$ 개 직선}에서 성립: $s=+4$ 최소벡터를 지나는 546개, $s=-4$ 의 546개, 그리고 $\sqrt3E_8$ 와 등거리인 546개(그중 91개가 위수 3 원소의 고정격자). (class number 를 쓰지 않는 증명: Hermite 상수 $\gamma_8=2$, $E_8$ 의 유일성, 7-design 성질.)
\end{enumerate}
\end{theorem}
증명의 핵심: $N(\nu)\in\{1,3\}$ 이면 기본단위 $\varepsilon=\tfrac{3+\sqrt{13}}2$ 의 거듭제곱을 곱해 노름 2 벡터가 생기고, $N(\nu)=4$ 는 $s=\pm6$ 인 최소벡터를 요구하는데 없다; $5,6,7,8$ 은 $O_F$ 의 노름이 아니다. 이득은 $g=2184/(2\cdot72)=15.17$.

\subsection{일반 $p$: 군론적 필요조건}
\begin{proposition}\label{prop:sl2pnec}\tagP\ (5회차; 심사 통과)\ $W_p$ 가 극대 order 위의 짝 unimodular $G$-격자라 하자.
\begin{enumerate}
\item (Cor.\ A$''$) $(p+1)/2$ 가 합성수이면 불안정: 비분리 토러스의 소수 위수 $q\mid\tfrac{p+1}2$ 원소($N=\tfrac{p+1}{2q}>1$)의 고정공간에서 $R_U\ge(\tfrac{Nq-1}{Nq-q})^{4(N-1)}$ ($p=17$: $9.99$; $29$: $11.8$; $41$: $12.5$). 심사자의 강화: unimodular 가 아니어도 $R_U(L)R_U(L^\#)>1$.
\item (Lemma B) 분리 토러스의 홀수 위수 부분군 $H$ 의 고정격자는 $W_p^H=N_H(W_p)$, $\det=|H|^{\dim W^H}$, \emph{정확히} $R_U=1$.
\item $(p+1)/2$ 가 소수이면($p=5,13,37,61,73,157,193,\dots$) 고정벡터를 갖는 부분군은 분리 토러스의 홀수 위수 부분군뿐이다.
\item (11회차 T2) 부분군 $H$ 의 $\Q$-isotypic 성분도 $\Delta$-부분공간이다: $SL_2(3)$ 의 Hurwitz 성분은 $R=2^{-e}\le1$; $p\equiv1\pmod5$ 에서 $SL_2(5)$ 의 icosian 성분은 $R=(625/256)^{2r/15}2^{-e}$ ($e$ = 2-진 접착 지수; $p=61$ 은 $e\ge3$ 이어야 안정).
\item (직선 기준) $R_U\le1\ (\forall\text{ 직선})\iff m_+(w)m_-(w)\ge\theta_p:=\tfrac{(p-1)^2}{4p}\ (\forall w\ne0)$ (주 직선의 경우; 일반적으로는 ideal 판), 여기서 $m_\pm$ 는 $\eta=\pm\sqrt p$ 고유공간 성분의 제곱노름. 동치로: $G$-불변 형식의 1-매개변수족 $Q_t=e^tm_++e^{-t}m_-$ (행렬식 1; $t$ 의 주기 $2\log\varepsilon$) 이 \emph{모든} $t$ 에서 최소값 $\ge(p-1)/\sqrt p$.
\end{enumerate}
\end{proposition}

\subsection{Mordell--Weil 재정식화 (11회차)}
\begin{theorem}\label{thm:MW}\tagP\ (본인 유도를 T1 에이전트가 검증·정정; Gross 1990 등의 원문은 열람하지 못함)\
$C_p:\ y^2=x^p-x$ (종수 $g=(p-1)/2$)에 $SL_2(p)$ 가 $x\mapsto\tfrac{ax+b}{cx+d}$, $y\mapsto y/(cx+d)^{(p+1)/2}$ 로 충실하게 작용하고 $H^1(C_p)=\psi+\psi'$. $\mathrm{Jac}(C_p)$ 는 superspecial 이고 Frobenius $\pi$ 는 $\pi^2=p$ 를 만족하며 $W$ 위에서 $\eta=\pm\pi$. $E$ 를 초특이 타원곡선이라 하면
\[
M_p:=\mathrm{Hom}(\mathrm{Jac}\,C_p,E)\quad(\text{랭크 }2(p-1),\ \text{형식 }2\deg)
\]
는 극대 order 를 갖는 $G$-격자이고, 짝 unimodular Weil 격자 $W_p$ 는 같은 $\Z[G]$-가군에 뒤틀린 형식 $\Tr_{F/\Q}(\tfrac{\varepsilon}{\sqrt p}\beta)$ 를 준 것이다. 형식의 원 $\{Q_t\}$ 위에서 unimodular 점과 Mordell--Weil 점($2\deg=\sqrt p\,Q$)은 정반대(거리 $\log\varepsilon$)에 있다. 직선 안정성은
\[
N_{F/\Q}\big(q_F(\phi)\big)\ \ge\ \Big(\frac{p-1}4\Big)^2\qquad(\forall\ \phi\ne0),\qquad q_F(\phi)=\tfrac12\big(\deg\phi+v(\phi)\sqrt p\big),\ \ v=\tfrac{\langle\phi\pi,\phi\rangle}{2p},
\]
와 동치이고, 특히 \emph{모든 비상수 사상 $C_p\to E$ 의 차수가 $\ge(p-1)/2=g$} 일 것을 요구한다.
\end{theorem}
\begin{itemize}
\item 초등적 한계(Castelnuovo--Severi/gonality)는 $\deg\ge(p+3)/4$ 뿐 --- 필요한 것의 절반.
\item $p\equiv5\pmod8$ 이면 $x\mapsto x^{(p-1)/4}$ 로 $C_p\to D:Y^2=X^5-X$, 이어서 $D\to E$ (차수 2): 차수 정확히 $(p-1)/2$ 의 사상 --- 이것이 분리 토러스 고정격자의 경계 벡터다($p=13$: $\sqrt3E_8\subset$ Leech). $p=13$ 에서는 차수 6 사상이 $6|G|$ 개 있고 그중 Galois 몫에서 오는 것은 $1/6$ 뿐이다.
\item 차수 $d$ 사상은 ``$R^*(e_1+\dots+e_4)\equiv\sum_{P\in\mathbb P^1(\mathbb F_p)}P\pmod2$ 인 차수 $d$ 유리함수 $R$'' 과 대응한다.
\item \tagH\ generic 격자에 대한 기대 위반 벡터 수(단위군 궤도당)는 $h\,\omega_h^2\theta_p^{h/2}\log\varepsilon$ ($h=p-1$): $1.2\times10^4$ ($p=13$), $1.6\times10^5$ ($37$), $246$ ($61$), $2.3$ ($73$), $2\times10^{-26}$ ($157$). $G$-궤도 수로는 약 $5,\ 3,\ 10^{-3},\ 6\times10^{-6}$. 즉 이 모형은 ``$p=37$ 은 불안정, $p\ge61$ 은 구조적 짧은 벡터가 없는 한 안정'' 을 예측한다(단 $p=13$ 에서는 과대예측: 실제 위반 0).
\end{itemize}

\subsection{직접 계산: $p=37$ 은 불안정 (10회차)}\label{sec:weilcomp}
홀함수 Weil 표현($T=\mathrm{diag}(\zeta^{x^2})$, $S=-(\zeta^{2xy}-\zeta^{-2xy})/\sqrt p$)을 $\Q(\zeta_p)$ 의 모든 매장으로 Minkowski 매장하고, 분리 토러스 $C_{(p-1)/4}$ 의 고정벡터 $v_0$ 로 $\Lambda=\Z[G]v_0$ 를 만든다. $B=\Tr_{K/\Q}\langle\cdot,\cdot\rangle$ 는 정수값이고 $D:=\det_B(\Lambda)^{1/(p-1)}=r^2p$, $\Gamma(w):=m_+(w)m_-(w)/D$ 로 두면 $\Gamma(w)<\theta_p\Rightarrow R_{\Delta w}\ge(\theta_p/\Gamma(w))^4>1$. 세 개의 독립 구현(하위 에이전트 2 + 저자의 프로토타입 \texttt{rounds/R10\_weil\_lattices/orchestrator/})이 일치했다.

\begin{center}
\small
\begin{tabular}{@{}rrrrlll@{}}
\toprule
$p$ & 랭크 & $D$ & $\theta_p$ & $\min\Gamma/\theta_p$ (찾은 값) & 판정 & 비고\\
\midrule
5 & 8 & 5 & 0.8 & 1 (전수) & 안정 & unimodular 점 $=E_8$\\
13 & 24 & 117 & 2.769 & 1 (전수; 세 류 각 6552쌍) & 정확히 경계 & unimodular 점 $=$ Leech\\
17 & 32 & & 3.765 & $<1$ & 불안정 & 비분리 위수 3 원소 (예측대로)\\
29 & 56 & & 6.759 & $45/49$; 고정 직선 $R_U=14.76$ & 불안정 & 비분리 위수 5, 3 원소 (예측대로)\\
\textbf{37} & 72 & 2997 & 8.757 & $\mathbf{7/9}$ & \textbf{불안정} & $R_U\ge(9/7)^4=2.7326$\\
\textbf{61} & 120 & & 14.754 & \S\ref{sec:p61} & \S\ref{sec:p61} & \\
\bottomrule
\end{tabular}
\end{center}

\begin{observation}[$p=37$]\label{obs:p37}\tagN\ (부동소수점 계산, 정수성 잔차 $\le2\times10^{-9}$; 독립 구현 일치)\
(i) $\Lambda$ 는 랭크 72, $D=2997=81\cdot37$; $t=\pm\log\varepsilon$ ($\varepsilon=6+\sqrt{37}$) 에서 $Q_t/\sqrt D$ 는 정수 형식 --- 72차원 짝 unimodular $SL_2(37)$-격자 $W_{37}$. (ii) 분리 토러스 $C_9$ 의 고정 직선은 $\log R_U=0.000000$, $\min\Gamma=\theta$ (보정). (iii) $C_3$-고정 부분격자(랭크 24, 전체로는 $R_U=1$)는 $\sqrt3\,\times$(근이 144개인 unimodular 랭크 24 격자 --- 즉 뿌리계가 $D_4^6$ 또는 $A_5^4D_4$ 인 Niemeier 격자). 그 근에 대응하는 $W_{37}$ 의 노름 6 벡터 72쌍은 모두 $s=\pm9$:
\[
N(\nu(w))=\frac{37\cdot6^2}{16}-\frac{9^2}4=63\ <\ 81=r^2 .
\]
(iv) 직접 인증: commutant order 의 작용으로 만든 랭크 8 부분격자 $O\cdot w\subset\Lambda\cap\Delta w$ 의 Gram 행렬식이 $((p/4)m_+m_-)^4$ 과 일치하고 $R_U\ge r^8D^4/\det(O\cdot w)=2.7326$. 이 하계는 $O$ 의 극대성을 가정하지 않으며, 같은 족의 모든 격자 $\Lambda I$ (모든 ideal 류)로 전이된다.
\end{observation}
따라서 $p=37$ 에서는 양자택일의 (b) 쪽이 성립한다: Weil 격자에 ``불균형한 짧은 벡터'' 가 있다. 구조적으로는 $p=13$ 의 $\sqrt3E_8$ (근의 $s$ 값이 허용 범위 안)이 $p=37$ 에서는 근을 가진 Niemeier 격자(근의 $s$ 값이 허용 범위 밖)로 바뀐 것이다. 이는 11회차 휴리스틱(위반 궤도 약 3개 예상; 발견된 것은 $703\times144=2|G|$ 개 벡터)과도 부합한다.

\subsection{$p=61$: 다음 사례 (10·12회차)}\label{sec:p61}
$p=61$: $\dim W=120$, $r=15$, $|G|=226920$, 이득 $g=63.0$, $\theta_{61}=14.754$, $\varepsilon=\tfrac{39+5\sqrt{61}}2$. unimodular 점은 120차원 짝 unimodular $SL_2(61)$-격자 $W_{61}$ 이고, 안정성은 (균형 잡힌 벡터에 대해) 최소노름 $\ge8$ 을 요구한다(노름 6 벡터는 $s$ 와 무관하게 위반; 노름 8 은 $|s|\le8$ 이어야 함).
\begin{observation}[10회차, 빌더 B]\tagN\
분리 토러스 $C_{15}$ 의 고정 직선(랭크 8), $C_5$-고정 부분격자(랭크 24), $C_3$-고정 부분격자(랭크 40)에서 $t$-원 전체(격자 간격 $0.1$)에 대해 반지름 $2\sqrt{\theta D}\cosh(0.05)$ 이하의 벡터를 전수 열거: \textbf{위반 없음}, 각각 $\min\Gamma=\theta_{61}$ 정확히(경계). 랭크 120 전체에 대해서는 $t$-격자 위의 LLL 기저와 그 쌍별 합에서 $\theta$ 미만이 없었다(약한 증거).
\end{observation}
즉 $p=37$ 을 무너뜨린 메커니즘($C_3$-고정 부분격자의 근)이 $p=61$ 에서는 나타나지 않는다. 12회차에서 부분군의 isotypic 성분($SL_2(5)$ 의 icosian 성분 등)에 대한 체계적 시험을 수행했다: 결과는 아래와 같다.


\subsection{조건부 주장의 올바른 형태 (11회차 T2 의 반박 시도 반영)}
\begin{proposition}[텐서족의 사영 분산]\tagP\ (Monte Carlo 확인)\
텐서족 $W_\Z\otimes_OL$ ($W$ 의 $\Delta$-랭크 $r$, $L$ 의 랭크 $k$)에서 전형적인 점의 사영 이차변동은 $q=q_0\big(1+\tfrac{r-1/2}{k}+O(\tfrac1{rk})\big)$, $q_0=1/\dim_\Q A$. 초과분 $r/k=n/(2D_k)$ 는 정확히 항등 방향(호흡 모드)의 기여 $1/n$ 이다.
\end{proposition}
따라서 기존 Klartag 회계를 그대로 쓰면 $V_{final}\approx V_0^{r/(k+r)}(|G|D_k/2e)^{k/(k+r)}$ 이고, \textbf{$k\approx r$ 에서는 지수 이득이 사라진다}($n^{2.0}$). $k=cr$ 이면 지수 $2+\tfrac{c-1}{2(c+1)}$, 완전한 이득은 $k\ge Cr\log n$ 에서만 나온다.
\begin{conjecture}[수정된 양자택일]\label{conj:dichotomy}\tagJ\ / \tagC\
$(p+1)/2$ 가 소수인 $p\equiv1\pmod4$ 에 대해 다음 중 하나가 무한히 자주 성립한다.
\begin{itemize}
\item[(a)] $W_p$ 가 모든 $\Delta$-부분공간에서 $R_U\le1$ 이고 사실상 경계인 부분공간들의 접촉 기여가 $o(D_k)$. 이 경우, $r\ge2$ 인 텐서족에 대한 등변 Klartag 논증(낮은 랭크 항과 Steinitz 류를 포함한 평균값 공식, 가중 궤도 계수; $k\ge Cr\log n$ 영역)이 증명된다면 $n=8rk\sim\tfrac C2p^2\log p$ 에서
\[
V_n\ \ge\ (1-o(1))\,\frac{p\,n^2}{2e}\ \ge\ c\,\frac{n^{5/2}}{\sqrt{\log n}}\qquad(\text{희소한 차원열}),
\]
즉 AGZ 가 전하는 Venkatesh 의 추측 $V_n\le n^2(\log n)^{O(1)}$ 의 반례.
\item[(b)] $W_p$ 는 텐서 불안정 --- $N_{F/\Q}(q_F(\phi))<((p-1)/4)^2$ 인 사상 $\phi:C_p\to E$ 가 존재($p=37$ 처럼).
\end{itemize}
\end{conjecture}
어느 쪽이든 증명되면 새로운 결과다. (a)에는 무한히 많은 그런 소수의 존재($q,2q-1$ 동시 소수 문제)도 필요하다. 반박 담당 에이전트는 (a)$\Rightarrow$결론의 경로에서 치명적 구조 장애를 찾지 못했다(평균값 공식의 주항 상수는 정확히 1; $O$-가군 제타 $\prod_{i<2r}\zeta_F(2k-i)$ 는 $k\ge r+1$ 에서 수렴하고 $\le1.17$; 강제되는 경계류가 흡수하는 제약 궤도는 $O(p)\ll D_k$). 그러나 theta 급수의 Hilbert 모듈러 형식 제약 등은 검토하지 못했고, 해석적 정리 자체가 미증명이므로 이 경로 전체는 \emph{추측}이다.
